\section{Model and Instrumentation Details}
\label{sec:system-details-app}

This appendix collects implementation details trimmed from the main paper for
page budget reasons.

\subsection{Imputation-Layer Architecture}

The imputation layer (Stage~1) is a compact \emph{permutation-invariant set encoder} over
the stale-quantile and feedback tokens. It is built so the same model consumes any
feedback-window size and is invariant to observation order:

\begin{enumerate}
  \item \textbf{Tokenization}: the normalized stale quantile boundaries, lightweight
        metadata, and up to $K$ recent feedback observations (predicate type, boundary,
        estimated and actual selectivity) are embedded into a common token space, with a
        validity mask for unused slots.
  \item \textbf{Set encoding}: the feedback tokens are aggregated by a permutation-invariant
        encoder (self-attention over the unordered window followed by masked pooling), so the
        representation is invariant to observation order and handles a variable number of
        observations.
  \item \textbf{Residual head}: the pooled representation, fused with the encoded prior,
        predicts a correction delta $\boldsymbol{\delta}$ to the interior quantile
        boundaries; the output is clamped and monotonized to a valid CDF, so the model learns
        drift-induced deviations rather than absolute boundaries.
\end{enumerate}

\paragraph{Training objective.}
Unlike a reconstruction-trained prior, the model is trained on the composite
optimizer-facing objective of the main paper: the held-out future-predicate log-Q-error of
the \emph{projected} marginal, a FactorJoin bin-mass join-regret term, an in-window
feedback-consistency residual, and a light reconstruction regularizer toward the post-drift
marginal. Training uses a sparse-to-dense feedback curriculum and a domain-randomized mixture
of drift families (batch loads, range/date shifts, skew evolution, outlier bursts, multimodal
and seasonal drift) so the completion generalizes across feedback budgets and drift
mechanisms. The model is trained offline in PyTorch with Adam and reused across columns
without per-table retuning. The main-paper ablation shows that applying the projection inside
the training loop is inessential---post-hoc projection performs identically---so the
contribution is the training objective, not in-loop calibration.

\subsection{Per-Conjunct Operator Instrumentation}

Filter and scan-filter-project operators are augmented with per-conjunct
selectivity counters. During execution, whenever a batch of rows is processed
through a compound predicate $\phi_1 \wedge \phi_2 \wedge \cdots \wedge \phi_m$,
each conjunct $\phi_j$ is evaluated independently on the full input batch:
\[
  n_j^{\text{in}} \mathrel{+}= |\text{batch}|, \qquad
  n_j^{\text{out}} \mathrel{+}= |\{r \in \text{batch} : \phi_j(r)\}|.
\]
The final intersection of all conjunct results produces the same output as
the original compound filter while additionally recording per-conjunct
independent selectivities $s^*_j = n_j^{\text{out}} / n_j^{\text{in}}$.
Upon operator completion, these counters are written to the engine's existing
runtime metrics store. In the deployment model used by OASIS, the collector
simply piggybacks on predicate evaluation already being performed by the
engine, flushes a structured record to disk, and emits at most one
observation tuple per filter conjunct. As a result, each query writes at
most as many observation records as it has conjunctive filter clauses.

\begin{algorithm}[H]\small
\caption{Per-Conjunct Operator Instrumentation}
\label{alg:instrument}
\begin{algorithmic}[1]
\Require Compound filter $\Phi = \phi_1 \wedge \cdots \wedge \phi_m$, input batch $B$
\Ensure  Filtered output $B'$, counters $n_j^{\text{in}}, n_j^{\text{out}}$
\State $\text{mask} \leftarrow [\textbf{true}]^{|B|}$
\For{$j = 1$ \textbf{to} $m$}
  \State $\text{result}_j \leftarrow \phi_j(B)$
  \State $n_j^{\text{in}} \mathrel{+}= |B|$; \quad $n_j^{\text{out}} \mathrel{+}= |\text{result}_j|$
  \State $\text{mask} \leftarrow \text{mask} \wedge \text{result}_j$
\EndFor
\State \Return $B[\text{mask}]$
\Statex
\Statex \textit{// On operator completion, emit per-conjunct metrics:}
\For{$j = 1$ \textbf{to} $m$}
  \State \Call{EmitMetric}{conjunct\_$j$, $n_j^{\text{in}}$, $n_j^{\text{out}}$}
\EndFor
\end{algorithmic}
\end{algorithm}

\subsection{QuickSel-H Adaptation Details}

The QuickSel system~\cite{park2020quicksel} is a selectivity learner that fits a
mixture-model density to satisfy feedback constraints. Because QuickSel outputs
point selectivity estimates rather than histogram boundaries, we evaluate a
clearly labeled \textbf{QuickSel-H} adaptation that bridges the gap between
selectivity learning and histogram correction. The adaptation performs three steps:

\begin{enumerate}
  \item \textbf{Mixture-model fitting}: Given the same stale prior $H_0$ and
        observation window $\mathcal{O}$ of size $K{=}16$, QuickSel-H fits a
        Gaussian mixture model with $M{=}5$ components whose CDF is constrained
        to match observed selectivities at the predicate values in $\mathcal{O}$.
        The optimization uses the same relative-entropy objective as the original
        QuickSel, solved via iterated scaling.
  \item \textbf{Quantile extraction}: From the fitted mixture CDF, we extract
        the $B{-}1{=}9$ quantile values at the equi-depth levels used by all other
        methods. This ensures QuickSel-H outputs on the same histogram grid.
  \item \textbf{Output interface}: The extracted quantiles are passed through
        the same monotonic projection as OASIS and returned as the corrected histogram.
\end{enumerate}

\noindent\textbf{Why this adaptation is necessary.}
QuickSel's native output is a density estimate over the full domain, not a
fixed-resolution histogram consumable by the CBO. To compare under a matched
output interface, we must extract quantiles from this density. The extraction
step introduces a minor information bottleneck (compressing a rich density into
$B{-}1$ boundaries), but this is identical to the bottleneck all other methods
face and does not fundamentally change QuickSel's correction mechanism. We do
not claim QuickSel-H is a line-by-line reproduction of the original system; it
is a faithful adaptation that preserves QuickSel's core density-fitting approach
while conforming to the common output format required by our evaluation protocol.
